\begin{explanationbox}
	\textbf{\textcolor{CExplanation}{一句话直觉}}

	假设你投篮 $n$ 次，每次命中概率 $p$，那么平均命中次数就是 $n$ 乘以 $p$，即 $E\xi=np$。

	\medskip
	\textbf{\textcolor{CExplanation}{核心拆解}}
	\begin{description}[style=nextline, leftmargin=2.8em, labelsep=0.5em, itemsep=0.45em]
		\item[\textbf{要点一（伯努利指示变量）：}]
		      定义 $X_i$ 表示第 $i$ 次试验是否成功：成功时 $X_i=1$，失败时 $X_i=0$。
		      则 $X_i\sim B(1,p)$，且
		      \[
			      EX_i=1\cdot p+0\cdot(1-p)=p.
		      \]

		\item[\textbf{要点二（期望可加）：}]
		      对于随机变量之和，只要各项期望存在，期望等于期望之和；
		      独立性不是期望可加的必要条件。
		      这里的独立重复条件主要用于保证 $\xi$ 服从二项分布。
		      因此
		      \[
			      E\xi
			      =
			      E\!\left(\sum_{i=1}^n X_i\right)
			      =
			      \sum_{i=1}^n EX_i
			      =
			      np.
		      \]
	\end{description}

	\medskip
	\textbf{\textcolor{CExplanation}{几何本质（分布重心）}}

	离散型随机变量的期望就是概率分布的重心，即所有可能取值以概率为权重的加权平均。二项分布的重心位于 $np$ 处。

	\medskip
	\textbf{\textcolor{CExplanation}{代数意义（加权平均）}}

	\[
		E\xi=\sum_{k=0}^n k\cdot P(\xi=k).
	\]
	代入二项分布概率公式后化简即为 $np$。因此它本质上是求和结果的简化。

	\medskip
	\textbf{\textcolor{CExplanation}{考点价值（直接应用）}}
	\begin{itemize}[leftmargin=*, itemsep=0.5em, topsep=0.4em]
		\item \textbf{考法一（求期望）：} 已知 $\xi\sim B(n,p)$，直接使用 $E\xi=np$ 求期望。
		\item \textbf{考法二（反求参数）：} 给出期望和 $n$ 或 $p$ 之一，反求另一个参数。
	\end{itemize}

	\medskip
	\textbf{\textcolor{CExplanation}{顿悟点}}

	\textbf{直接套二项分布公式前要先判断模型；单独求成功次数期望时，更本质的方法是指示变量分解与期望线性性。}

	\medskip
	\textbf{\textcolor{CExplanation}{使用场景}}
	\begin{itemize}[leftmargin=*, itemsep=0.5em, topsep=0.4em]
		\item \textbf{场景一（直接计算）：} 已知二项分布模型，直接使用 $np$ 求期望。
		\item \textbf{场景二（线性变形）：} 对于 $Y=a\xi+b$，利用线性性质 $EY=aE\xi+b$，将 $a,b$ 与 $np$ 结合。
		\item \textbf{场景三（非二项但求期望）：} 若 $\xi$ 是若干指示变量之和，可用 $E\xi=\sum EX_i$ 求期望，不必强行判断二项分布。
	\end{itemize}
\end{explanationbox}